\cotefile{15}
\scanpage{34}

% B1 | text | head and first five lines
\uline{Ex. 3} Soit fixé, \unsure{un|une} \unsure{fibre|filtre} de la catégorie des
\unsure{Dessin|Dessins}$_{\otimes}$ $C$ sur $A$ comme dans Ex. 2, avec
élément $\tau$ de $V_A(C)$. Si $E$ est \inserted{\unsure{centrale|centré}} de rang $n$,
il faut considérer que les $C$ \uline{généralisent} \inserted{les \deleted{(\unsure{de même si on se restreint})}}
\unsure{des|les} liens $\unsure{G|G_m}$ \uuline{sur $A$} trivialisés par \unsure{$K$|un},

% B2 | text | fine writing between lines five and six
(qui correspond au cas où \deleted{\unsure{sur}}, dans \unsure{tt}, ordre maximal est \deleted{\unsure{un}}
\deleted{et \unsure{donc} (cas où $E$ est \unsure{déployée} sur \unsure{tt} complété de \unsure{$A$})}
une algèbre d'Azumaya sur $R$, i.e. la classe de \unsure{lien} $\mathrm{Br}(K)$ est dans
l'image de $\mathrm{Br}(R)$)

% B3 | text | left margin opposite opening
% slanted
\margin{gerbe sur $A$ de lien $G_m$ \inserted{et trivialisée par \unsure{$K$}} $\longmapsto$ catégorie $C$ des représentations de cette gerbe qui sont de degré $n$, poids $1$}

% B4 | diagram | left margin below the first marginal note
% slanted
\margin{\diagram{\(\begin{array}{c@{\quad}l}
 C \longmapsto & \text{gerbe des}\\[-2pt]
 \text{\scriptsize cat. de} & \text{foncteurs}\\
 \text{\scriptsize \unsure{Dessin}}_{\otimes} & \text{admissibles}\\
 \text{\scriptsize \unsure{munie} $\tau$} & \text{de $C$ dans les}\\
 \text{\scriptsize où $\{\tau\}\in$} & \text{modules \unsure{sur}}\\
 \operatorname{Im}(\mathrm{Br}\,A\to\mathrm{Br}\,K) & \text{(trivialisée \unsure{par}}\\
 & \text{\unsure{de} $C$)}
\end{array}\)}}

% B5 | text | continuation through middle of sheet
et il \deleted{faut} voir de \unsure{$=$|même} que les couples $(C,\tau)$
généralisent la donnée d'une gerbe de groupe $\mu_n$.

% B6 | text | middle, new paragraph
Pour ceci, \deleted{\unsure{on}} notons \deleted{qu'\unsure{une} \unsure{gerbe} \inserted{la catégorie des}} qu'une
gerbe de lien $\mu_n$ sur $A$ définit bien une $(C,\tau)$ :
$C$ : les \deleted{\unsure{trivialisations} de} représentations de degré $n$, de rang $1$
de la gerbe, $\tau=$ \deleted{la représentation déterminant}

% B7 | text | middle, new paragraph beginning Nous
Nous allons déterminer d'autre part les autoéquivalences de $(\unsure{E|C},\tau)$ \unsure{;}
en termes d'un objet $X$ de $C$, \inserted{d'où $O=\operatorname{End}(X)$, $\nu$}
c'est la catégorie des \deleted{\unsure{objets}} bimodules
$P$ sous $O$, inv. loc. sur $A$ des deux cotés,
munis d'un isom \inserted{dans $\mathbb P(A)$} ($\deleted{\nu(X)=L\otimes}\;\tau=L\otimes_A\nu(X)$)

% B8 | display | lower middle, centred norm
\[
 N_{O/A}(P_O)\simeq\deleted{\unsure{O}} A.
\]

% B9 | text | left margin of norm paragraph, braced
\margin{$\varphi(X)=P\otimes_O X$\newline
$\nu(\varphi(X))=N(P)\otimes\nu(X)$\newline
$\phantom{\nu(\varphi(X))}\unsure{\wr\!|\,|\,\simeq}$\newline
$\varphi(\nu(X))$}

% B10 | text | below norm and outlined cancellation
\deleted{compatible avec}.
\deleted{Supposons que tous les \unsure{$P$} \inserted{de \unsure{$\mathbb P(O/A)$}} soient isomorphes
ce qui est le cas si $A$ est local, ou semi local. Alors}

% B11 | text | foot of sheet
C'est une \unsure{Gr.|Gp.}-catégorie \inserted{et $=$ une catégorie de Picard strt. \uline{comm.}}, où les endom.
de l'objet $1$ sont les éléments $\lambda\in A^*$ tels que $\lambda^n=1$ ;
\deleted{\unsure{cette}} catégorie de Picard contient comme sous-catégorie pleine $\underline{\mathrm{Tors}}(S,\mu_n)$ des
% archive sheet number: 33
